% deleted_paragraph.tex
% Archive of material removed from Draft6.tex to keep the manuscript clean.
% NOTHING here is compiled into the paper (everything sits inside \iffalse ... \fi).
% Parts: 1) deleted original introduction  2) superseded phrases  3) reviewer comments.
% The pristine pre-cleanup source is preserved in Draft6.tex.orig.

\iffalse

================================================================
 PART 1.  DELETED ORIGINAL INTRODUCTION (superseded by current intro)
================================================================

% --- PARA 1: hook. ---
Whether liquid water is a single continuum of local structures \comment{[Be careful with word choice. It is definitely a single substance. The question is whether there are two structural states within the one substance.]} or a mixture of two distinct structural states has been debated since the nineteenth century \cite{whiting-1883, rontgen-1892} \comment{[If you want to talk about this from the nineteenth century, find citations that discussed this from the nineteenth century and cite them here]}, and the debate has proved remarkably hard to settle. This question matters because water is anomalous: it reaches its maximum density at $4\,^\circ$C, its
isothermal compressibility rises on supercooling, and its dynamics cross
over from fragile to strong ~\cite{gallo_water_2016}. \comment{[The next sentence finally goes to ``two pictures''. But it has already been logcally detached from the previous sentence.]} The microscopic picture of water needs to explain these behaviors, and two pictures compete to do so. Continuum models treat water as a broad, unimodal distribution of local environments~\cite{eisenberg_structure_2005}, whereas mixture models\cite{whiting-1883,rontgen-1892,wernet-2004, huang-2009, russo_understanding_2014} \comment{[Why here? With the two dashes in front and after, it is a ``word meander''. Also, this is not the only work on mixture models. Citing only this does not capture the big picture of the field.]} posit two or more distinct populations of local molecular environment \comment{[populations of what?]}: a low-density, tetrahedrally ordered form, and a denser, less-ordered form, whose relative fractions shift with temperature and pressure. The two pictures have proved remarkably hard to separate because they reproduce the same macroscopic behavior: decades of thermodynamic and scattering data are consistent with both, so unambiguous evidence for either has remained elusive. \cite{gallo_water_2016, handle_supercooled_2017}. 

\comment{[Narten's work is not about macroscopic behavior or any of the thermodynamic anomalies. It is X-ray diffraction on the structure. Are you sure this is correctly cited and it is what you want ot say? I think it is not consistent with the sentence written here. Eisenberg is a whole book -- is the whole book consistent with the sentence before this?]}
\comment{[The logic from sentence to sentence and within each sentence is so vague in this paragraph.]}

% --- PARA 2: explain two state model--
Recent computational work by Shi and Tanaka has given the two-state picture considerably firmer footing~\cite{shi_direct_2020,russo_understanding_2014,shi_microscopic_2018}. In this framework, which we refer to hereafter as the two-state model, water is a dynamic mixture of locally favored tetrahedral structures (LFTS)---stabilized by four hydrogen bonds, with low density and high local symmetry---and disordered normal-liquid structures (DNLS), with broken tetrahedral symmetry, higher density, and greater entropy. Central to their analysis is the
structural descriptor $\zeta$, which quantifies translational order in the second coordination shell~\cite{russo_understanding_2014}. The distribution $P(\zeta)$ is bimodal in several classical water models, and the
oxygen--oxygen partial structure factor $S(k)$ carries two overlapping peaks within its apparent first diffraction peak: a first sharp diffraction peak at $k_{T1} = k r_{\mathrm{OO}}/2\pi \approx 3/4$ from the tetrahedral
LFTS geometry, and an ordinary-liquid peak at $k_{D1} \approx 1$ from DNLS~\cite{shi_direct_2020,shi-2019}. These two peak positions are a concrete, first-principles prediction of the model---and the target any independent
test of the two-state picture must reproduce. \comment{[No. Two-peak in a distribution function does not imply there is a two-state picture. ]}

% --- PARA 3: ML lit review + gap, merged. Lands circularity as the trap. ---
Machine learning offers a way to test this prediction without presupposing it. \comment{[You need to provide some background before making the assertion. Additionally, this sounds like a sweeping statement, as machine learning is way too general.]} Neural networks can identify phases and locate transitions directly from raw configurations, with no hand-tuned order parameters~\cite{carrasquilla_machine_2017}, and water has become a natural proving ground \comment{[Why is water ``natural proving ground''? What are we trying to prove?]}. Gartner and co-workers paired a near-quantum machine-learned potential with free-energy sampling and found behavior consistent with a liquid--liquid transition~\cite{gartner_signatures_2020}; Donkor and co-workers showed that only nonlocal descriptors resolve high- and low-density domains near the critical point~\cite{donkor_beyond_2024}; and Li and co-workers trained an unsupervised autoencoder on large molecular-dynamics trajectories and recovered molecular-level evidence for two interconvertible structures~\cite{li_evidence_2026}. \comment{[The previous few sentences were exactly like the example I talked about in the presentation: ``Dad is mowing, mom is shopping, I am listening to music.'' Where is the connection and logic?]} However, a common limitation runs through the structural evidence for two states. \comment{[A sweeping statement again.]} Earlier structure-factor work \comment{[What's the ``earlier structure-factor'' work that you are referring to?]} resolved the two populations largely through a single hand-crafted descriptor and a $\zeta$-binned scattering decomposition, so the classification and its reciprocal-space signature were read off the same quantity~\cite{shi_direct_2020}; the more recent deep-learning approaches recovered the states from opaque, tuned reaction coordinates validated only against an internal volume--fraction relation~\cite{li_evidence_2026} \comment{[Is \cite{li_evidence_2026} the ``more recent deep-learning appraoches''? What are they?]}. In each case the evidence is, to some degree,
circular: the observable used to \emph{define} the states is essentially the one used to \emph{confirm} them. What has been missing is a classification validated against a physical observable the classifier never sees. \comment{[You might be saying something here. But it is not clear to me what you mean. Try harder? If this is an important statement, it needs to be a paragraph in itself.]}
\comment{[When you discuss other people's work, you need to be very specific.]}

\comment{[After you have made the scientific gap clear, introduce what you are doing in this work.]}

================================================================
 PART 2.  SUPERSEDED SENTENCES / PHRASES (earlier \delete markup)
================================================================

- hybrid density-denoising pipelines yield the clearest

- Liquid water is one substance, but whether its molecules occupy a continuous distribution of local environments or fluctuate between two distinguishable structural states has been debated since the nineteenth century.

- R{\"o}ntgen proposed a two-state interpretation of liquid water in 1892, while subsequent structural studies and modern simulations have continued to ask whether a tetrahedrally ordered environment can be distinguished from a more densely packed, disordered environment~\cite{rontgen-1892,whiting-1883,narten_observed_1969}.

- This question is not merely terminological. Any microscopic description must account for water's unusual macroscopic behavior, including its density maximum near $4\,^{\circ}$C, the increase of its isothermal compressibility on supercooling, and its crossover from fragile to strong dynamics~\cite{gallo_water_2016}.

- The difficulty is that these anomalies can be described either by a continuous distribution of local structures or by a temperature- and pressure-dependent mixture of local environments.

- In the latter picture, the relevant populations are populations of molecules: one population has relatively low density and high tetrahedral order, whereas the other has higher density, weaker tetrahedral order, and greater configurational entropy.

- Experimental scattering, structural analysis, and thermodynamic measurements have provided evidence relevant to both descriptions, so the macroscopic anomalies alone do not identify which microscopic picture is correct.

- is

- hybrid

- molecules in the low-density transition zone
    flagged as noise are shown in grey

- The silhouette score for this
    pipeline is $0.2841$, the highest among the methods tested.

- The signature also survives a change of water model. Running the same pipeline on three models (TIP4P/2005, TIP5P, and SWM4-NDP) at $T=-20\,^\circ$C (Fig.~\ref{fig:generality}d), we find a clear first sharp diffraction peak in every LFTS water group (solid curves)---tetrahedral ordering is a robust feature of the LFTS population, strongest in TIP5P, then TIP4P/2005, then SWM4-NDP. The DNLS group (dashed) are more uniform, their ordinary-liquid peak at $k_{D1} \approx 1.05$ holding steady across models, with SWM4-NDP showing the strongest DNLS signal.

- clustering

- (silhouette score $0.2841$) by discarding the ambiguous transition-zone molecules before classification

- This two-state signature is not specific to a single state point. Across the supercooled-to-ambient range ($-30$ to $+30\,^\circ$C), the first sharp diffraction peak strengthens monotonically upon cooling as the LFTS fraction grows toward the low-density amorphous limit, while the DNLS peak remains essentially stable; and the signature persists across three independent water models (TIP4P/2005, TIP5P, and SWM4-NDP), with tetrahedral ordering of the LFTS population a robust feature in every case. The convergence of ML clustering in order-parameter space and structure-factor validation in reciprocal space therefore provides a compelling, model-independent case for genuine structural heterogeneity in liquid water.

- clustering

- clustering-plus-validation

- test the data

- across a hard boundary

================================================================
 PART 3.  REVIEWER COMMENTS moved out of the main text
   (status OPEN unless a resolution was recorded this round)
================================================================

[OPEN]
   context: ...the geometry of each water molecule. The relative density of feature points projected on to the $q$-$\zeta$ plane is shown in  Fig.~\ref{fig:classification_comparison}a,
   comment: [Figure 2 cannoty be mentioned before Figure 1. Please change figure order. This has to be Figure 1a.]

[ADDRESSED]
   context: ...in the region where the modes overlap. We therefore allow the classification to leave such molecules unassigned rather than force every molecule  into an assigned group.
   comment: [The logic does not follow after this sentence. We have talked about this earlier --- you don't know whether your classified groups are LFTS and DNLS until \emph{you check the $\zeta$ parameter and structure factor}. Please smoothen the overall logical flow.]
   note:    [ADDRESSED 2026-07-23: the two groups are now called high-/low-$\zeta$ throughout the method; the names LFTS and DNLS are introduced only at the structure-factor section, once the tetrahedral vs. close-packed peaks are shown.]

[OPEN]
   context: ...on between confidently assigned and genuinely ambiguous molecules, and setting some molecules aside would amount to imposing an arbitrary cut through a populated region.
   comment: [This sentence does not make sense because water has two structural states but also a continuum. How is there any ``distinction between confidently assigned and genuinely ambiguous molecules'' in this case?]

[ADDRESSED]
   context: ...kelihood ratio \begin{equation} \Lambda(\mathbf{x}) = \log p_{\mathrm{LFTS}}(\mathbf{x}) - \log p_{\mathrm{DNLS}}(\mathbf{x}), \label{eq:likelihood_ratio} \end{equation}
   comment: [How are $ p_{\mathrm{LFTS}}(\mathbf{x})$ and $ p_{\mathrm{DNLS}}(\mathbf{x})$ defined or calculated? ]
   note:    [ADDRESSED 2026-07-23: text now states they are the two fitted GMM component densities, $p_{high}$ and $p_{low}$, identified by which component has the larger/smaller mean $\zeta$.]

[OPEN]
   context: ...ahedral state (LFTS) or the disordered normal-liquid state (DNLS); molecules for which neither state is clearly favored ($|\Lambda|<\tau$) are set aside as transitional.
   comment: [Based on the LFTS graph and the transitional graph -- DNLS should be a water holding a hexagon rather than a mixture of stuff? Also, there is no reason that DNLS should look smaller than LFTS.]

[ADDRESSED]
   context: ...assigns the remainder to LFTS ($32.5\%$ of all molecules) and DNLS ($57.7\%$). The resulting partition has the highest silhouette score ($0.34$) of the methods we tested
   comment: [will get this value t]
   note:    [ADDRESSED: likelihood-ratio silhouette $=0.34$ (four-feature); K-means $0.2418$ and GMM $0.2144$ retained.]

[OPEN]
   context: ...sparse transitional band.  However, neither the silhouette score nor these structural signatures establish that the split is physically real. Both are internal measures.
   comment: [What does ``internal meastures'' mean?]

[ADDRESSED]
   context: ...e two dense lobes and the sparse     interconverting region between them.  The silhouette score for this     pipeline is $0.34$, the highest among the methods tested
   comment: [update this sihouette score]
   note:    [ADDRESSED: $0.34$.]

[ADDRESSED]
   context: ...nimum, confirming that the two populations are physically distinct     rather than an artifact of the classification scheme.   }   \label{fig:validation} \end{figure*}
   comment: Generality. What do your results contribute to the debate?
   note:    [ADDRESSED: the temperature-dependence paragraphs below now answer this---two states at low $T$ that homogenize on warming; carried into the Conclusions.]

\fi
